\documentclass[12pt]{article}
\usepackage[utf8]{inputenc}
\usepackage[margin=1in]{geometry}
\usepackage{times}
\usepackage{setspace}
\usepackage{graphicx}
\usepackage{listings}
\usepackage{xcolor}
\usepackage{booktabs}
\usepackage{amsmath}
\usepackage{amsfonts}

\setstretch{1.5}

\title{A Hybrid Framework for Demand Forecasting: Integrating Domain-Specific Context with Statistical Machine Learning}
\author{Sidra Shaikh \\ Department of Management Information Systems \\ University of Houston}
\date{June 1, 2026}

\begin{document}
\maketitle

\begin{abstract}
Predicting retail demand accurately is a major challenge for supply chain planning. Traditional forecasting tools usually look only at past sales and pricing numbers, which means they tend to struggle when sudden real-world shifts happen, such as bad weather or public health crises. In this project, I built a hybrid demand forecasting framework that combines historical sales data with external context scores for weather and epidemic conditions. Using a Random Forest regressor trained on 76,000 transaction records, I compared a standard baseline model against my context-aware hybrid model. The results show that the baseline model produced a Mean Absolute Percentage Error (MAPE) of 43.51\%, while the hybrid model lowered that error to 39.37\%, achieving a 4.14 percentage point improvement. This project demonstrates that simple, structured feature engineering can significantly improve prediction accuracy without requiring complex deep learning systems.
\end{abstract}

\section{Introduction}
In retail and supply chain management, having accurate demand forecasts is essential for keeping inventory balanced and avoiding stockouts. Most standard forecasting methods rely strictly on past sales records and pricing. However, when external disruptions occur—like severe weather storms or health restrictions—these historical models often fail because they assume past trends will simply repeat themselves.

In this project, I wanted to test whether adding practical context features into a machine learning model could make retail demand predictions more reliable. Rather than using complex deep learning architectures, I focused on a clean, interpretable approach: a Random Forest regressor combined with custom numeric weights for weather conditions and epidemic statuses. By mapping qualitative categories into simple scaling scores, I aimed to give the model better visibility into real-world factors that drive consumer behavior.

The main contributions of this work include:
\begin{enumerate}
    \item Designing a simple data pipeline that converts categorical weather and epidemic data into numerical weighting scores.
    \item Implementing and evaluating a Random Forest model on 76,000 retail transaction records, comparing a baseline feature set against the hybrid context-aware feature set.
    \item Using a modified MAPE evaluation metric that adds an epsilon smoothing factor ($\epsilon = 1.0$) to the denominator, preventing errors from blowing up on rows with zero sales.
\end{enumerate}

\section{Related Work}
Demand forecasting has evolved from simple moving averages and classical time series models, such as ARIMA (Box et al. \cite{box2015}; Hyndman and Athanasopoulos \cite{hyndman2018}), to machine learning regression techniques. Breiman \cite{breiman2001} established the foundation for Random Forest regression, showing how ensemble tree methods handle non-linear interactions well without overfitting. 

In supply chain analytics, researchers have increasingly looked at hybrid modeling strategies (Zhang \cite{zhang2003}) to combine traditional baselines with external indicators. Recent industry benchmarks, such as the M4 Competition findings summarized by Makridakis et al. \cite{makridakis2020}, highlight that incorporating external context helps bridge the stability gap in retail forecasting. My work builds on these ideas by testing a lightweight feature engineering approach on a large retail dataset.

\section{Data Dictionary}
The dataset used in this research comes from Kaggle and contains exactly 76,000 transaction records. It covers 5 stores and 20 products over a 760-day period from January 1, 2022, to January 30, 2024, with no missing values. Table 1 outlines the primary features used in the models.

\begin{table}[h]
\centering
\caption{Variables and Definitions}
\begin{tabular}{lll}
\toprule
Feature Name & Data Type & Description \\
\midrule
Units Sold & Numerical & Total units sold per day (Target Variable) \\
Inventory Level & Numerical & Current stock available \\
Units Ordered & Numerical & Quantity ordered for replenishment \\
Price & Numerical & Retail product price \\
Discount & Numerical & Applied discount percentage \\
Competitor Pricing & Numerical & Competitor pricing benchmark \\
Weather Condition & Categorical & Local weather state (Sunny, Cloudy, Rainy, Snowy) \\
Epidemic & Binary & Public health crisis indicator flag (0 or 1) \\
\bottomrule
\end{tabular}
\end{table}

\section{Methodology}
The project workflow is divided into data preprocessing, feature engineering, and model training.

\subsection{Feature Engineering}
To help the machine learning model understand external conditions, I converted categorical weather and epidemic labels into numerical scores using a dictionary lookup in `src/features/create_hybrid_features.py`:
\begin{itemize}
    \item \textbf{Weather Score}: Sunny ($1.2$), Cloudy ($1.0$), Rainy ($0.9$), Snowy ($0.8$).
    \item \textbf{Epidemic Score}: Normal ($1.0$), Flagged outbreak ($0.6$).
\end{itemize}

\subsection{Model Training and Evaluation}
I trained two separate Random Forest regression models using Python's \textit{scikit-learn} library (`n_estimators=100`, `random_state=42`), using an 80/20 train/test split (60,800 training rows and 15,200 testing rows):
\begin{itemize}
    \item \textbf{Baseline Model} (`src/models/train_baseline.py`): Trained using `['Inventory Level', 'Units Ordered', 'Price', 'Discount', 'Competitor Pricing']`.
    \item \textbf{Hybrid Model} (`src/models/train_hybrid.py`): Trained using `['Inventory Level', 'Units Ordered', 'Price', 'Weather_Score', 'Epidemic_Score']`.
\end{itemize}

Because the dataset includes transactions with zero units sold, standard MAPE calculations would result in division-by-zero errors. To fix this, I used a modified MAPE formula with an epsilon smoothing factor ($\epsilon = 1.0$) added to the denominator:
\begin{equation}
\text{MAPE} = \frac{1}{n} \sum_{i=1}^{n} \left| \frac{y_i - \hat{y}_i}{y_i + 1} \right| \times 100
\end{equation}

\begin{lstlisting}[language=Python, caption=Hybrid Model Implementation Sample, frame=single]
import pandas as pd
from sklearn.ensemble import RandomForestRegressor
from sklearn.model_selection import train_test_split

data = pd.read_csv('../data/processed/demand_forecasting_hybrid.csv')
features = ['Inventory Level', 'Units Ordered', 'Price', 'Weather_Score', 'Epidemic_Score']
X = data[features]
y = data['Units Sold']

X_train, X_test, y_train, y_test = train_test_split(
    X, y, test_size=0.2, random_state=42
)

model = RandomForestRegressor(n_estimators=100, random_state=42)
model.fit(X_train, y_train)
predictions = model.predict(X_test)
\end{lstlisting}

\section{Results and Discussion}
When evaluated on the test set, the baseline model produced a MAPE of 43.51\%. The hybrid model achieved a MAPE of 39.37\%, representing a 4.14 percentage point improvement (about a 9.5\% relative error reduction). 

During periods of bad weather or epidemic disruptions, the baseline model tended to overpredict sales because it relied only on historical sales trends. The hybrid model adjusted its predictions downward more appropriately, helping prevent excess inventory buildup.

\subsection{Limitations}
There are several limitations to keep in mind:
\begin{enumerate}
    \item \textbf{Feature Trade-off}: The hybrid model swapped out `Discount` and `Competitor Pricing` for the context scores, so the accuracy improvement cannot be credited solely to the weather and epidemic scores by themselves.
    \item \textbf{Random Split}: The data split was done randomly rather than chronologically by date, which means there could be some temporal overlap between training and testing sets.
    \item \textbf{Manual Weights}: The weather and epidemic scores were assigned manually based on intuition rather than being learned automatically from the data.
\end{enumerate}

\section{Conclusion}
This research shows that incorporating simple context scores into a Random Forest model can improve retail demand forecasting accuracy. By adding weather and epidemic multipliers, the hybrid model reduced prediction error from 43.51\% to 39.37\%. For future work, I plan to explore automating these score assignments using real-time data feeds.

\newpage
\bibliographystyle{plain}
\bibliography{references}

\end{document}